% Auto-generated by notebooks/03_evaluate.ipynb cell case-study-cell
% Не редактировать вручную — перезапустить ячейку.
% Источник: datasets/processed/case_study_traceback.json

\subparagraph*{Кейс 1. Регулярное выражение как ложный сигнал (filled-pasta на сыре)}

\noindent
\textbf{Тип кейса:} Ошибочный. \quad
\textbf{Код товара:} \texttt{00154222}. \quad
\textbf{Категория:} pasta.
\quad\textbf{Класс ошибки:} regex-false-positive.

\smallskip

\noindent\textbf{Входные данные от партнёра:}
\begin{itemize}
  \item \emph{product\_name}: Three Cheese Ravioli
  \item \emph{brands}: M\&S Food, Marks \& Spencer
  \item \emph{quantity}: 400 g
  \item \emph{ingredients\_text}: Tomatoes (39\%), Cooked Egg Pasta (24\%) (Durum Wheat Semolina, Water, Pasteurised Egg, Wheatflour), Three Cheese Mix (23\%) (Ricotta Cheese (Milk), Extra Mature…
\end{itemize}

\noindent\textbf{Прохождение каскада по атрибутам:}

{\small
\begin{xltabular}{\textwidth}{|X|c|c|c|c|c|c|c|c|c|}
\caption{Прохождение каскада для товара \texttt{00154222}.} \\ \hline
Атрибут & Сл.\,1 & Сл.\,2 & Conf. & Порог & Сл.\,3 & Сл.\,4 & Выбран & Ответ & Эталон \\ \hline
\endfirsthead
\multicolumn{10}{l}{\textit{Продолжение таблицы}} \\ \hline
Атрибут & Сл.\,1 & Сл.\,2 & Conf. & Порог & Сл.\,3 & Сл.\,4 & Выбран & Ответ & Эталон \\ \hline
\endhead
is\_filled & false (h) & true & 0,97 & 0,90 & — & — & Прав.h & false & true $\times$ \\ \hline
\end{xltabular}
}

\noindent\textbf{Объяснение.} ML слой корректно распознаёт начинку (ml\_pred=true, ml\_conf=0.97 > порога 0.9). Однако высокоприоритетное эвристическое правило по продукту <<Three Cheese Ravioli>> форсирует значение false (предположение, что любое блюдо со словом cheese — это пицца или сэндвич, а не паста с начинкой). Поскольку правила high-tier перекрывают решение ML, итоговый ответ системы — false, что не совпадает с эталоном (agreement\_ratio=1.0). Этот пример демонстрирует, что класс regex-false-positive состоит преимущественно из случаев, когда правило подавляет правильный ML-ответ. Вывод (зафиксирован в INBOX): правило по cheese на is\_filled следует ослабить или ограничить контекстом (исключить Ravioli, Tortellini, Cappelletti).

\medskip

\subparagraph*{Кейс 2. Смешение близких классов (cream vs fresh)}

\noindent
\textbf{Тип кейса:} Ошибочный. \quad
\textbf{Код товара:} \texttt{3184670017166}. \quad
\textbf{Категория:} cheeses.
\quad\textbf{Класс ошибки:} class-confusion.

\smallskip

\noindent\textbf{Входные данные от партнёра:}
\begin{itemize}
  \item \emph{product\_name}: Chevre a tartiner
  \item \emph{brands}: Rians
  \item \emph{quantity}: 100 g
  \item \emph{ingredients\_text}: Lait de chèvre, protéine de lait de chèvre, crème pasteurisée de chèvre, sel, ciboulette (0,4\%), ferments lactiques
\end{itemize}

\noindent\textbf{Прохождение каскада по атрибутам:}

{\small
\begin{xltabular}{\textwidth}{|X|c|c|c|c|c|c|c|c|c|}
\caption{Прохождение каскада для товара \texttt{3184670017166}.} \\ \hline
Атрибут & Сл.\,1 & Сл.\,2 & Conf. & Порог & Сл.\,3 & Сл.\,4 & Выбран & Ответ & Эталон \\ \hline
\endfirsthead
\multicolumn{10}{l}{\textit{Продолжение таблицы}} \\ \hline
Атрибут & Сл.\,1 & Сл.\,2 & Conf. & Порог & Сл.\,3 & Сл.\,4 & Выбран & Ответ & Эталон \\ \hline
\endhead
texture & — & fresh & 0,87 & 0,65 & — & — & ML & fresh & cream $\times$ \\ \hline
\end{xltabular}
}

\noindent\textbf{Объяснение.} Слой ML с высокой уверенностью выбирает класс fresh для мягкого козьего сыра <<Chevre a tartiner>>. Эталон же относит товар к крем-сырам (cream). Граница между fresh-cheese и cream-cheese размыта на уровне OFF-тегов и проявляется как смешение классов в матрице ошибок (§3.3.x). Такая ошибка ожидаема и потенциально снимается экспертной заменой границ в схеме (см. INBOX, ticket по reschema cheeses.texture).

\medskip

\subparagraph*{Кейс 3. Silver-шум: товар не относится к категории}

\noindent
\textbf{Тип кейса:} Ошибочный. \quad
\textbf{Код товара:} \texttt{2000000074755}. \quad
\textbf{Категория:} pasta.
\quad\textbf{Класс ошибки:} silver-noise.

\smallskip

\noindent\textbf{Входные данные от партнёра:}
\begin{itemize}
  \item \emph{product\_name}: Pâté de foie de volaille
  \item \emph{brands}: AJM, Sarl AJM
  \item \emph{quantity}: 136 gr
  \item \emph{ingredients\_text}: 70\% de gras de porc ; 30\% de foie de volaille ; Sel ; Poivre ; Muscade
\end{itemize}

\noindent\textbf{Прохождение каскада по атрибутам:}

{\small
\begin{xltabular}{\textwidth}{|X|c|c|c|c|c|c|c|c|c|}
\caption{Прохождение каскада для товара \texttt{2000000074755}.} \\ \hline
Атрибут & Сл.\,1 & Сл.\,2 & Conf. & Порог & Сл.\,3 & Сл.\,4 & Выбран & Ответ & Эталон \\ \hline
\endfirsthead
\multicolumn{10}{l}{\textit{Продолжение таблицы}} \\ \hline
Атрибут & Сл.\,1 & Сл.\,2 & Conf. & Порог & Сл.\,3 & Сл.\,4 & Выбран & Ответ & Эталон \\ \hline
\endhead
is\_gluten\_free & — & false & 0,98 & 0,90 & — & — & ML & false & true $\times$ \\ \hline
\end{xltabular}
}

\noindent\textbf{Объяснение.} Товар <<Pâté de foie de volaille>> попал в выборку pasta из-за тегов OFF, однако к пасте отношения не имеет. Silver-эталон присвоил is\_gluten\_free=true, тогда как каскад абстрагируется (cascade\_pred=false по умолчанию из-за низкой уверенности ML). Эталон в данном случае некорректен (пшеничная мука обычно встречается в подобных паштетах, поэтому is\_gluten\_free=true сомнительно). Каскад уверенно выдаёт is\_gluten\_free=false (ml\_conf=0.98). Низкое значение consensus\_agreement=0,67 подтверждает спорность эталона. Пример иллюстрирует ограничение метода получения silver-эталона из OFF-тегов и обосновывает необходимость пересчёта метрик на human-gold консенсусе (§3.3.4).

\medskip

\subparagraph*{Кейс 4. Сложный мультиязычный товар, ML слой даёт уверенный ответ по всем атрибутам}

\noindent
\textbf{Тип кейса:} Успешный. \quad
\textbf{Код товара:} \texttt{20114992}. \quad
\textbf{Категория:} pasta.

\smallskip

\noindent\textbf{Входные данные от партнёра:}
\begin{itemize}
  \item \emph{product\_name}: Bio Spaghetti
  \item \emph{brands}: Combino
  \item \emph{quantity}: 500 g
  \item \emph{ingredients\_text}: Hartweizengriess
\end{itemize}

\noindent\textbf{Прохождение каскада по атрибутам:}

{\small
\begin{xltabular}{\textwidth}{|X|c|c|c|c|c|c|c|c|c|}
\caption{Прохождение каскада для товара \texttt{20114992}.} \\ \hline
Атрибут & Сл.\,1 & Сл.\,2 & Conf. & Порог & Сл.\,3 & Сл.\,4 & Выбран & Ответ & Эталон \\ \hline
\endfirsthead
\multicolumn{10}{l}{\textit{Продолжение таблицы}} \\ \hline
Атрибут & Сл.\,1 & Сл.\,2 & Conf. & Порог & Сл.\,3 & Сл.\,4 & Выбран & Ответ & Эталон \\ \hline
\endhead
grain\_type & — & wheat & 0,96 & 0,90 & — & — & ML & wheat & wheat \checkmark \\ \hline
is\_filled & false (h) & false & 1,00 & 0,90 & — & — & Прав.h & false & false \checkmark \\ \hline
is\_organic & — & true & 0,98 & 0,85 & — & — & ML & true & true \checkmark \\ \hline
is\_gluten\_free & — & false & 0,99 & 0,90 & — & — & ML & false & false \checkmark \\ \hline
pasta\_shape & — & spaghetti & 0,86 & 0,80 & — & — & ML & spaghetti & spaghetti \checkmark \\ \hline
is\_vegan & — & true & 0,98 & 0,70 & — & — & ML & true & true \checkmark \\ \hline
cuisine\_origin & — & italian & 0,94 & 0,60 & — & — & ML & italian & italian \checkmark \\ \hline
\end{xltabular}
}

\noindent\textbf{Объяснение.} Товар <<Bio Spaghetti>> бренда Combino — мультиязычный (немецкое Bio + итальянское Spaghetti). Регулярное выражение распознаёт is\_filled=false по правилу высокого приоритета (нет слов начинки). ML слой уверенно предсказывает все остальные шесть атрибутов (cuisine\_origin=italian, grain\_type=wheat, is\_organic=true, is\_vegan=true, pasta\_shape=spaghetti, is\_gluten\_free=false) с уверенностью выше соответствующих порогов. Эскалация в LLM не требуется. Пример демонстрирует корректность многоязычных эмбеддингов (paraphrase-multilingual-MiniLM-L12-v2) и адекватность per-attribute порогов.

\medskip

\subparagraph*{Кейс 5. Французская тёмная плитка с начинкой — комбинированная работа правил и ML}

\noindent
\textbf{Тип кейса:} Успешный. \quad
\textbf{Код товара:} \texttt{3250392007737}. \quad
\textbf{Категория:} chocolate.

\smallskip

\noindent\textbf{Входные данные от партнёра:}
\begin{itemize}
  \item \emph{product\_name}: Tablette fourrée chocolat au lait caramel à la fleur de sel de Guérande - 125g
  \item \emph{brands}: Ivoria
  \item \emph{quantity}: 125 g
  \item \emph{ingredients\_text}: Sucre, caramel au lait frais d'Isigny et à la fleur de sel de Guérande (30\%) [sirop de glucose, lait frais d'Isigny (9\%), sucre, eau, beurre concentré (dont la…
\end{itemize}

\noindent\textbf{Прохождение каскада по атрибутам:}

{\small
\begin{xltabular}{\textwidth}{|X|c|c|c|c|c|c|c|c|c|}
\caption{Прохождение каскада для товара \texttt{3250392007737}.} \\ \hline
Атрибут & Сл.\,1 & Сл.\,2 & Conf. & Порог & Сл.\,3 & Сл.\,4 & Выбран & Ответ & Эталон \\ \hline
\endfirsthead
\multicolumn{10}{l}{\textit{Продолжение таблицы}} \\ \hline
Атрибут & Сл.\,1 & Сл.\,2 & Conf. & Порог & Сл.\,3 & Сл.\,4 & Выбран & Ответ & Эталон \\ \hline
\endhead
chocolate\_type & — & milk & 0,97 & 0,90 & — & — & ML & milk & milk \checkmark \\ \hline
contains\_nuts & false (h) & false & 0,97 & 0,70 & — & — & Прав.h & false & false \checkmark \\ \hline
is\_filled & true (h) & true & 0,70 & 0,55 & — & — & Прав.h & true & true \checkmark \\ \hline
chocolate\_extra & with\_caramel (l) & with\_caramel & 0,91 & 0,85 & — & — & ML & with\_caramel & with\_caramel \checkmark \\ \hline
is\_organic & — & false & 0,99 & 0,90 & — & — & ML & false & false \checkmark \\ \hline
flavor\_profile & — & salty\_caramel & 0,88 & 0,50 & — & — & ML & salty\_caramel & salty\_caramel \checkmark \\ \hline
\end{xltabular}
}

\noindent\textbf{Объяснение.} Шоколад <<Tablette fourrée chocolat au lait caramel à la fleur de sel de Guérande>> бренда Ivoria — тёмная плитка с начинкой солёная карамель. Правила высокого приоритета корректно срабатывают на is\_filled=true (по слову fourrée) и contains\_nuts=false. Правило низкого приоритета корректно фиксирует chocolate\_extra=with\_caramel. Оставшиеся три атрибута (chocolate\_type=milk, is\_organic=false, flavor\_profile=salty\_caramel) корректно предсказывает ML слой. Эскалация в LLM не требуется. Пример иллюстрирует кооперативную работу слоёв на товаре со сложным составом.

\medskip
